% Example content adapted from David Molik’s tenure document.
\PortfolioSection{Summary of Candidate’s Directed Service}
\SectionGuidance{Describe work performed as part of your assigned responsibilities. Explain the quality and impact of that work, your individual contributions, and the evidence supporting your account.}
\ExampleNote

\subsection*{Ability to Serve the Public}

As Department Head for the Center for Scholarly Publishing (now the Open Publishing Exchange), I served the Kansas State University community by advancing library services that directly support faculty, researchers, and students. A significant component of this work involved engaging with campus research leadership to better understand needs related to scholarly publishing, data management, and compliance.

In April 2025, I presented the scope, landscape, and strategic direction of the Center to the Associate Deans for Research Council. This engagement strengthened working relationships with research administrators and positioned the Libraries as an active partner in the University’s research enterprise. As a result of these discussions, library publishing and data services were more clearly aligned with faculty workflows, funder requirements, and institutional research priorities.

Through my leadership of the Open Publishing Exchange, I also supported student success and public access to scholarship by expanding and stabilizing open publishing and open educational resource services, ensuring that faculty and students have access to affordable, high-quality instructional and research materials.

\subsection*{Evidence of Specialization}

I demonstrated professional specialization in scholarly communication, publishing infrastructure, and research data governance through the development and coordination of new and evolving library services. I authored a comprehensive departmental landscape document and executive summary that articulated the scope, priorities, and strategic direction of scholarly publishing and data services at Kansas State University. These documents now serve as foundational references for service planning, decision-making, and communication with campus partners.

I also oversaw early development of a Research Data Management and Curation program aligned with FAIR principles and federal open-access mandates, including the Nelson Memo. This work required specialized knowledge of data standards, repository ecosystems, and research compliance and involved coordination with the Office of Research and Research IT to ensure alignment with campus infrastructure planning, including CRIS/RIMS discussions.

In addition, I led the reorganization of governance structures for the Open and Alternative Textbook Initiative (OATI), drafting bylaws for a revised committee structure that clarified roles, responsibilities, and decision-making authority. This work reflects specialization in open education policy, publishing workflows, and sustainable service models.

\subsection*{Supervisory and Administrative Expertise}

Throughout the evaluation period, I exercised supervisory and administrative expertise as Department Head, providing leadership and oversight for a team of Scholarly Communication librarians and associated services. I balanced strategic planning with day-to-day operational management, ensuring continuity of services while advancing new initiatives.

I demonstrated administrative effectiveness by setting departmental priorities, coordinating across library units, and representing the department in campus-wide planning conversations. My participation in leadership development programs, including Harvard’s Leadership Institute for Academic Librarians and the Kansas State University Fall New Department Head Institute, directly informed my approach to personnel management, budgeting, communication, and change management.

\Redacted{Directed Service Evidence}{Internal planning documents, private consultation
updates, and confidential evidence links are omitted. Attach appropriate evidence
in your private submission.}
